%%% The main file. It contains definitions of basic parameters and includes all other parts.

%% Settings for single-side (simplex) printing
% Margins: left 40mm, right 25mm, top and bottom 25mm
% (but beware, LaTeX adds 1in implicitly)
\documentclass[12pt,a4paper]{report}
\setlength\textwidth{145mm}
\setlength\textheight{247mm}
\setlength\oddsidemargin{15mm}
\setlength\evensidemargin{15mm}
\setlength\topmargin{0mm}
\setlength\headsep{0mm}
\setlength\headheight{0mm}
% \openright makes the following text appear on a right-hand page
\let\openright=\clearpage

%% Settings for two-sided (duplex) printing
% \documentclass[12pt,a4paper,twoside,openright]{report}
% \setlength\textwidth{145mm}
% \setlength\textheight{247mm}
% \setlength\oddsidemargin{14.2mm}
% \setlength\evensidemargin{0mm}
% \setlength\topmargin{0mm}
% \setlength\headsep{0mm}
% \setlength\headheight{0mm}
% \let\openright=\cleardoublepage

%% Character encoding: usually latin2, cp1250 or utf8:
\usepackage[utf8]{inputenc}

%% Further useful packages (included in most LaTeX distributions)
\usepackage{amsmath}        % extensions for typesetting of math
\usepackage{amsfonts}       % math fonts
\usepackage{amsthm}         % theorems, definitions, etc.
\usepackage{bbding}         % various symbols (squares, asterisks, scissors, ...)
\usepackage{bm}             % boldface symbols (\bm)
\usepackage{graphicx}       % embedding of pictures
\usepackage{fancyvrb}       % improved verbatim environment
\usepackage[numbers]{natbib}     
\bibliographystyle{unsrtnat} % order citations by their appearance
\usepackage[nottoc]{tocbibind} % makes sure that bibliography and the lists
			    % of figures/tables are included in the table
			    % of contents
\usepackage{dcolumn}        % improved alignment of table columns
\usepackage{booktabs}       % improved horizontal lines in tables
\usepackage{paralist}       % improved enumerate and itemize
\usepackage[usenames]{xcolor}  % typesetting in color
\usepackage[sc]{mathpazo} % different font to make it a little bit more fun

%%% Basic information on the thesis

% Thesis title in English (exactly as in the formal assignment)
\def\ThesisTitle{Data-driven Web Application Generator}

% Author of the thesis
\def\ThesisAuthor{Bc. Tobiáš Potoček}

% Year when the thesis is submitted
\def\YearSubmitted{2016}

% Name of the department or institute, where the work was officially assigned
% (according to the Organizational Structure of MFF UK in English,
% or a full name of a department outside MFF)
\def\Department{Department of Software Engineering}

% Is it a department (katedra), or an institute (ústav)?
\def\DeptType{Department}

% Thesis supervisor: name, surname and titles
\def\Supervisor{RNDr. Jiří Helmich}

% Supervisor's department (again according to Organizational structure of MFF)
\def\SupervisorsDepartment{Department of Software Engineering}

% Study programme and specialization
\def\StudyProgramme{Department of Software Engineering}
\def\StudyBranch{ITI}
% * <tobiaspotocek@gmail.com> 2016-06-07T14:17:16.541Z:
%
% > ITI
%
% Is this correct?
%
% ^.

% An optional dedication: you can thank whomever you wish (your supervisor,
% consultant, a person who lent the software, etc.)
\def\Dedication{%
Dedication.
}

% Abstract (recommended length around 80-200 words; this is not a copy of your thesis assignment!)
\def\Abstract{%
Abstract.
}

% 3 to 5 keywords (recommended), each enclosed in curly braces
\def\Keywords{%
{Linked Data} {LinkedPipes} {visualization} {application} {generator}
}
% * <tobiaspotocek@gmail.com> 2016-06-07T14:20:02.620Z:
%
% > {Linked Data} {LinkedPipes} {visualization} {application} {generator}
%
% This is going to need sume extra work.
%
% ^.

%% The hyperref package for clickable links in PDF and also for storing
%% metadata to PDF (including the table of contents).
\usepackage[pdftex,unicode]{hyperref}   % Must follow all other packages
\hypersetup{breaklinks=true}
\hypersetup{pdftitle={\ThesisTitle}}
\hypersetup{pdfauthor={\ThesisAuthor}}
\hypersetup{pdfkeywords=\Keywords}
\hypersetup{urlcolor=blue}

% Definitions of macros (see description inside)
%%% This file contains definitions of various useful macros and environments %%%
%%% Please add more macros here instead of cluttering other files with them. %%%

%%% Minor tweaks of style

% These macros employ a little dirty trick to convince LaTeX to typeset
% chapter headings sanely, without lots of empty space above them.
% Feel free to ignore.
\makeatletter
\def\@makechapterhead#1{
  {\parindent \z@ \raggedright \normalfont
   \Huge\bfseries \thechapter. #1
   \par\nobreak
   \vskip 20\p@
}}
\def\@makeschapterhead#1{
  {\parindent \z@ \raggedright \normalfont
   \Huge\bfseries #1
   \par\nobreak
   \vskip 20\p@
}}
\makeatother

% This macro defines a chapter, which is not numbered, but is included
% in the table of contents.
\def\chapwithtoc#1{
\chapter*{#1}
\addcontentsline{toc}{chapter}{#1}
}

% Draw black "slugs" whenever a line overflows, so that we can spot it easily.
\overfullrule=1mm

%%% Macros for definitions, theorems, claims, examples, ... (requires amsthm package)

\theoremstyle{plain}
\newtheorem{thm}{Theorem}
\newtheorem{lemma}[thm]{Lemma}
\newtheorem{claim}[thm]{Claim}

\theoremstyle{plain}
\newtheorem{defn}{Definition}

\theoremstyle{remark}
\newtheorem*{cor}{Corollary}
\newtheorem*{rem}{Remark}
\newtheorem*{example}{Example}

%%% An environment for proofs

%%% FIXME %%% \newenvironment{proof}{
%%% FIXME %%%   \par\medskip\noindent
%%% FIXME %%%   \textit{Proof}.
%%% FIXME %%% }{
%%% FIXME %%% \newline
%%% FIXME %%% \rightline{$\square$}  % or \SquareCastShadowBottomRight from bbding package
%%% FIXME %%% }

%%% An environment for typesetting of program code and input/output
%%% of programs. (Requires the fancyvrb package -- fancy verbatim.)

%%% By tobice: we override the original verbatim environment as 
%%% the verbatim environment is highlighted with a different color
%%% in Overleaf (doesn't work for code)
\DefineVerbatimEnvironment{verbatim}{Verbatim}{fontsize=\scriptsize, frame=single}

%%% The field of all real and natural numbers
\newcommand{\R}{\mathbb{R}}
\newcommand{\N}{\mathbb{N}}

%%% Useful operators for statistics and probability
\DeclareMathOperator{\pr}{\textsf{P}}
\DeclareMathOperator{\E}{\textsf{E}\,}
\DeclareMathOperator{\var}{\textrm{var}}
\DeclareMathOperator{\sd}{\textrm{sd}}

%%% Transposition of a vector/matrix
\newcommand{\T}[1]{#1^\top}

%%% Various math goodies
\newcommand{\goto}{\rightarrow}
\newcommand{\gotop}{\stackrel{P}{\longrightarrow}}
\newcommand{\maon}[1]{o(n^{#1})}
\newcommand{\abs}[1]{\left|{#1}\right|}
\newcommand{\dint}{\int_0^\tau\!\!\int_0^\tau}
\newcommand{\isqr}[1]{\frac{1}{\sqrt{#1}}}

%%% Various table goodies
\newcommand{\pulrad}[1]{\raisebox{1.5ex}[0pt]{#1}}
\newcommand{\mc}[1]{\multicolumn{1}{c}{#1}}


% Title page and various mandatory informational pages
\begin{document}
%%% Title page of the thesis and other mandatory pages

%%% Title page of the thesis

\pagestyle{empty}
\hypersetup{pageanchor=false}
\begin{center}

\centerline{\mbox{\includegraphics[width=166mm]{img/logo-en.pdf}}}

\vspace{-8mm}
\vfill

{\bf\Large MASTER THESIS}

\vfill

{\LARGE\ThesisAuthor}

\vspace{15mm}

{\LARGE\bfseries\ThesisTitle}

\vfill

\Department

\vfill

\begin{tabular}{rl}

Supervisor of the master thesis: & \Supervisor \\
\noalign{\vspace{2mm}}
Study programme: & \StudyProgramme \\
\noalign{\vspace{2mm}}
Study branch: & \StudyBranch \\
\end{tabular}

\vfill

% Zde doplňte rok
Prague \YearSubmitted

\end{center}

\newpage

%%% Here should be a bound sheet included -- a signed copy of the "master
%%% thesis assignment". This assignment is NOT a part of the electronic
%%% version of the thesis. DO NOT SCAN.

%%% A page with a solemn declaration to the master thesis

\openright
\hypersetup{pageanchor=true}
\pagestyle{plain}
\pagenumbering{roman}
\vglue 0pt plus 1fill

\noindent
I declare that I carried out this master thesis independently, and only with the cited
sources, literature and other professional sources.

\medskip\noindent
I understand that my work relates to the rights and obligations under the Act No.~121/2000 Sb.,
the Copyright Act, as amended, in particular the fact that the Charles
University in Prague has the right to conclude a license agreement on the use of this
work as a school work pursuant to Section 60 subsection 1 of the Copyright Act.

\vspace{10mm}

\hbox{\hbox to 0.5\hsize{%
In ........ date ............	% FIXME!
\hss}\hbox to 0.5\hsize{%
signature of the author
\hss}}

\vspace{20mm}
\newpage

%%% Mandatory information page of the thesis

\openright

\vbox to 0.5\vsize{
\setlength\parindent{0mm}
\setlength\parskip{5mm}

Title:
\ThesisTitle

Author:
\ThesisAuthor

\DeptType:
\Department

Supervisor:
\Supervisor, \SupervisorsDepartment

Abstract:
\Abstract

Keywords:
\Keywords

\vss}

\newpage

%%% Dedication

\openright

\noindent
\Dedication

\newpage

\openright
\pagestyle{plain}
\pagenumbering{arabic}
\setcounter{page}{1}


%%% A page with automatically generated table of contents of the master thesis

\tableofcontents

%%% Each chapter is kept in a separate file
\chapter*{Introduction}
\addcontentsline{toc}{chapter}{Introduction}

The Internet of today is full of data and the amount is constantly increasing. By data we mean literary anything that can be found online, may it be a huge source of organized information like Wikipedia or a small table created with Microsoft Excel by a local government office. Unfortunately, our ability to use the whole potential of all currently available data is very limited. We want to to explore the data, we want to connect and combine different data sets, we want to look for patterns in the data, we want to make the data available to end users. All of that is possible, yet usually time and resources demanding. 

Let us consider for example the aforementioned "Excel table".  XLSX is a proprietary format that can be officially opened only by a paid proprietary software. This software offers a large yet still limited set of tools. It may very well happen that it will not allow us to do what we want. Given the topic of this thesis, let us say that we want to create from the table a new specialized interactive visualization (or an \textit{application}) and publish it online. To get this done,  we would most likely need a computer specialist or  even a developer. Those people are relatively rare and expensive.

Tools that attempt to solve this particular problem are called application generators. They allow non-developers to do developers' work which in this case means to generate from the data an interactive (possibly online) application. Such generators already exist, for example for the aforementioned XLSX format (or in general for tabular data) but they are usually very limited. One reason for this limitation is that the tabular data carry very little semantic information with them. One row represents an entity and each column contains one entity attribute. The first row might (or might not) contain labels for the attributes. That is all that the tabular data based generator can work with.

One of the positive recent trends is that more and more information is published using the Resource Description Framework \cite{rdf} while utilizing the Linked Data model \cite{ld}. Besides being open and offering a natural way of linking and combining different data sets, this framework is able to express any kind of information by extending the raw data with machine-readable semantic meta-data. In other words, data represented in RDF are, unlike the tabular data, not limited in what semantic information they can carry with them. A Linked Data driven application generator could benefit significantly from this feature. As such a generator could literary \textit{understand} the data, it would allow the user to generate way more versatile applications, tailored for all different kinds of data. 
% * <tobiaspotocek@gmail.com> 2016-06-09T10:02:42.253Z:
%
% >  A Linked Data driven application generator could benefit significantly from this feature
%
% Perhaps incorporate somewhere that there already exist lots of existing Linked Data applications but the goal is to make the creation process of new ones simpler.
%
% ^ <tobiaspotocek@gmail.com> 2016-06-09T10:10:59.089Z:
%
% Eg. There is an application called Justinian. It's hard to imagine that a non-developer user could create such a complex application. But the framework could be of great help.
%
% ^.
š
The main goal of this thesis will be to create such a Linked Data driven application generator. That means, as already suggested, that we will be focusing on allowing the non-developers do developers' work. Obviously, that is possible only to a certain extent. The data represented in RDF can be understood with software, but first someone (the developer) has to teach the software how to understand it. For example, for data containing geospatial information (RDF offers a way how to describe such information), the developer would have to extend the generator with a plugin showing such data on a~map.

\section*{Motivation and contribution}
% * <tobiaspotocek@gmail.com> 2016-06-08T21:14:06.686Z:
%
% > Motivation and contribution}
%
% This part should probably go to the Analysis chapter and the Introduction should be kept more general.
%
% ^.
Let us describe more in depth how our application generator will work and how it will help its users. To make it more organized, we will separate the descriptions by different type of users that should benefit from our generator. For each type of user, we will describe his role, his motivation (or perhaps a use case typical for this user) and eventually the contribution of this generator, i. e., how this generator is going to help him.

\begin{itemize}
\item \textbf{Developer}

\textit{Description:} A software engineer with deep technical knowledge.

\textit{Motivation/Use case:} A developer is given a task at his job to create an interactive online map application that visualizes provided company geospatial data. The developer is smart and wants to avoid doing the same thing again the next time. Therefore, he would prefer to create a universal solution that would allow other non-developer employees to create the applications on their own and do it with any data containing geospatial information. Also, it is very likely that he will be given similar task in the near future, but with a different data type and the developer would like to re-use as much work as possible from the first time.

\textit{Contribution:} Thanks to our application generator, the developer will have to focus only on developing the actual visualizer plugin. The generator will work both as a \textit{framework} and a \textit{platform}, solving many problems up front. As as a \textit{platform}, the generator will handle user agenda, registration of new plugins, application management (creating, updating, deleting, publishing), common application configuration and so on. As a \textit{framework}, the generator will define a standard  interface which will allow seamless integration of new plugins. Plus it will provide many drop-in solutions for typical challenges that the developer might encounter across multiple plugins (e.g. caching, multi-language support, support for fixing missing or incorrect data labels, re-usable user interface components etc.).

\item \textbf{Data Analyst}

\textit{Description:} A data analyst is not necessarily a developer, but he has enough knowledge to work with data on a lower level, typically convert data to RDF format or to augment existing RDF data sets with different meta-data. 

\textit{Motivation/Use case:} Let us say that a data analyst works at the same company as our developer. He has a couple of interesting data sets that he would like to make available for other employees so that they could use them to generate applications with the plugin developed by the developer.

\textit{Contribution:} By implementing the interface defined by the framework, each visualizer plugin describes what the input data should look like. Therefore the data analyst has clear instructions on how to prepare a data set to make it compatible with the visualizer plugin. After preparing such data set, the data analyst will be able to upload it to the generator and make it available for all other users that can use the data set for generating new applications.

\item \textbf{User}

\textit{Description: }This is the non-developer with no technical knowledge. He does possess at least some basic understanding of the data though.

\textit{Motivation/Use case: }A user wants to create a new interactive application from a data set and make it available online. For example, he might want to create a map application from the data set with geospatial information, but would be interested only in certain type of data from a certain area.

\textit{Contribution: }The application generator will let this user leverage the work that has been already done by the developer and the data analyst. He will be able to select a prepared data set (or choose his own), the generator will analyze the content of the data set and based on the results, the generator will offer the user a list of visualization plugins that are compatible with the data set. By choosing one plugin, the user will proceed to a configuration phase where he will be allowed to create the application according to his needs. The configuration possibilities will differ depending on the selected visualizer plugin, but typically the user will be allowed filter the data (select a subset) and to tune the level of interactivity for the audience. For example, he could either create a completely static visualization with pre-filtered data, or he could let the audience decide what they want to see. This will greatly increase the re-usability of data sets and visualizer plugins. A single data source with a single visualization plugin will work as a potential source for many applications, each serving a different purpose. None of these applications will involve any more work from the developer or from the data analyst.

Before publishing the application, the user will be allowed to fill missing information (e.g. missing labels in target language) and provide basic application meta data (name and description). After publishing, the application will be available online from a unique URL.

\item \textbf{Audience}

\textit{Description:} The end-user with no knowledge at all, just some basic computer skills.

\textit{Motivation/Use case:} Audience wants to use the application and wants to understand it.

\textit{Contribution:} The application generator will greatly simplify the process of creating new applications which will just by itself result in an increased number of generated applications. Therefore the audience will benefit just from the bare existence of the generator because without the generator, there would be way less applications and the data potential would be wasted. Plus thanks to the configuration phase, the applications will not be just raw visualizations and they will offer better comfort and user experience for the audience.
\end{itemize}

\section*{Goal of this thesis}

The goal of this thesis will be to create a Linked Data based application generator with all the features listed in the previous section. The application generator will work as a \textit{framework}, easily extendable with new visualization plugins that will allow generating applications from different types of data. It will define a clear interface that the new plugins will have to implement. Also many documented drop-in solutions for typical problems and situations will be available for the developer.

The application generator will also be a \textit{platform}, providing space for developers to register new plugins, space for data analyst to publish new data sets and space for common users to generate, manage and publish their applications.

Besides the generator itself, we will also implement two visualizer plugins to demonstrate the abilities of the \textit{framework} and the \textit{platform}.

\section*{Structure of the text}

In Chapter 1, we briefly describe RDF and Linked Data and provide related information necessary for understanding the rest of the thesis. In Chapter 2, we focus on similar existing tools and compare them to our proposed generator. In Chapter 3, we examine in depth the tool LinkedPipes as we will base our generator on top of it.  In Chapter 4, we look more in detail into the required features and we present steps and decisions necessary for achieving our goals. In this chapter, we also analyze and defend our decision to use LinkedPipes as a base for our generator, plus we describe how it has to be modified to fit our intentions. In Chapter 5, we dive into implementation details. Especially we focus on the framework aspect. We show how a new visualizer plugin can be integrated into the generator and what tools are available for the developer. In Chapter 6, we describe in detail the two visualizer plugins, D3.js Chord Visualizer and Google Maps Visualizer, that we implemented to demonstrate our generator's abilities. Related to the D3.js Chord Visualizer, there is a short theoretical part in this chapter that deals with representing and sampling graph structures (actual graphs as understood in graph theory) using RDF and SPARQL queries. In the final Chapter 7, we describe possible future work.
\chapter{Title of the first chapter}

An~example citation: \cite{Andel07}


\section{Title of the first subchapter of the first chapter}

\section{Title of the second subchapter of the first chapter}

\chapter{Title of the second chapter}

\section{Title of the first subchapter of the second chapter}

\section{Title of the second subchapter of the second chapter}

\chapter*{Conclusion}
\addcontentsline{toc}{chapter}{Conclusion}

We implemented the proposed \emph{application generator} which enables users to generate interactive web applications from RDF data. The user starts by selecting a source of RDF data. The \emph{application generator} analyzes the data and offers a list of \emph{visualizers} that are able to visualize the data. The user selects the \emph{visualizer} he is interested in and uses it to generate a new application. The \emph{application generator} allows the user to configure the application and then publish it.

We examined several existing related tools which helped us to specify the list of features that we decided our \emph{application generator} should support. This also showed that compared to other data representation formats, Linked Data offer the biggest potential for automatic analysis and automatic application generation.

Our \emph{application generator} is integrated into LinkedPipes Visualization. It utilizes its underlying LDVM implementation for fully automated and easily extendable Linked Data analysis. We extended LinkedPipes Visualization with all the necessary functionality required by our \emph{application generator}. Specifically, we added support for two new RDF vocabularies, RGML and FRESNEL.

The biggest asset of \emph{application generator} is the \emph{configuration phase}. Each \emph{visualizer} provides a \emph{configurator} user interface unique to the type of data the \emph{visualizer} supports. The \emph{configurator} interface enables the user to configure the application before it gets published. The configuration possibilities differ depending on the \emph{visualizer}, but typically the user is allowed to filter the data and tune the level of the application interactivity for the audience. It is also possible for example to fix missing text labels for individual entities in the data set.

Our \emph{application generator} is \emph{non-developers friendly}. The whole process of generating a new application requires zero technical knowledge and anyone with the base computer skills is able to generate a new application.

Our \emph{application generator} works as a \emph{platform}. Every user has its own account which he uses to manage his published application. The \emph{application generator} contains a catalog of published applications and also a simple repository of available data source that any user can utilize to generate applications.

Our \emph{application generator} is a \emph{framework}. Firstly, the \emph{framework} defines a clear way how new \emph{visualizer} can be seamlessly integrated into the \emph{application generator} which makes it easily extendable. We explained the process of integration a new \emph{visualizer} in a step-by-step guide which is part of this thesis. Secondly, the \emph{framework} provides the developer with several read-to-use drop-in solutions for some common tasks that can significantly speed up the process of developing a new \emph{visualizer}. Those common tasks for example include saving and loading the application configuration or multiple language support.

To showcase the capabilities of the \emph{platform} and the \emph{framework}, we developed two \emph{visualizers}. The first one is D3.js Chord Visualizer which visualizes graph data represented in RDF with a chord diagram. While making this \emph{visualizer}, we implemented and analyzed algorithms for fulltext search, for calculating a graph adjacency matrix and for graph sampling, all of which over a graph represented in RDF. The second one is Google Maps visualizer which visualizes geospatial data on a map. This \emph{visualizer} already existed in LinkedPipes Visualization and we decided to re-implement it for our \emph{application generation} to show the benefits of the \emph{application generator} concept over LinkedPipes Visualization.

Clearly, our \emph{application generator} is only as powerful as its \emph{visualizers}. At the current state, there are only two very domain specific \emph{visualizers} which means that the generator will yield no results for most of the RDF data. It is important to note that unlike for example tabular data, the RDF data are very versatile in what kind of information they can carry. That makes it very hard (possibly impossible) to create an almighty universal \emph{visualizer} capable of visualizing any RDF data (unless the \emph{visualizer} is very simplistic). However, we are confident that given the way the \emph{application generator} is designed, being both a \emph{platform} and a \emph{framework}, it lays a solid ground for future potential extensions. We believe we fulfilled the goal of this thesis.



%%% Bibliography
%%% Bibliography (literature used as a source)
%%%
%%% We employ bibTeX to construct the bibliography. It processes
%%% citations in the text (e.g., the \cite{...} macro) and looks up
%%% relevant entries in the bibliography.bib file.
%%%
%%% The \bibliographystyle command selects, which style will be used
%%% for references from the text. The argument in curly brackets is
%%% the name of the corresponding style file (*.bst). Both styles
%%% mentioned in this template are included in LaTeX distributions.



\renewcommand{\bibname}{Bibliography}

%%% Generate the bibliography. Beware that if you cited no works,
%%% the empty list will be omitted completely.

\bibliography{bibliography}

%%% If case you prefer to write the bibliography manually (without bibTeX),
%%% you can use the following. Please follow the ISO 690 standard and
%%% citation conventions of your field of research.

% \begin{thebibliography}{99}
%
% \bibitem{lamport94}
%   {\sc Lamport,} Leslie.
%   \emph{\LaTeX: A Document Preparation System}.
%   2nd edition.
%   Massachusetts: Addison Wesley, 1994.
%   ISBN 0-201-52983-1.
%
% \end{thebibliography}


%%% Figures used in the thesis (consider if this is needed)
\listoffigures

%%% Tables used in the thesis (consider if this is needed)
%%% In mathematical theses, it could be better to move the list of tables to the beginning of the thesis.
\listoftables

%%% Abbreviations used in the thesis, if any, including their explanation
%%% In mathematical theses, it could be better to move the list of abbreviations to the beginning of the thesis.
\chapwithtoc{List of Abbreviations}

%%% Attachments to the master thesis, if any. Each attachment must be
%%% referred to at least once from the text of the thesis. Attachments
%%% are numbered.
%%%
%%% The printed version should preferably contain attachments, which can be
%%% read (additional tables and charts, supplementary text, examples of
%%% program output, etc.). The electronic version is more suited for attachments
%%% which will likely be used in an electronic form rather than read (program
%%% source code, data files, interactive charts, etc.). Electronic attachments
%%% should be uploaded to SIS and optionally also included in the thesis on a~CD/DVD.
\chapwithtoc{Attachments}

\openright
\end{document}
